% Paper 2, §12 — Adaptation and merging
% Anchors: kb/notes/training/adaptation-and-merging.md
%          kb/excerpts/{hu2021-lora,dettmers2023-qlora,ilharco2022-task-vectors,
%                       wortsman2022-model-soups,yadav2023-ties-merging}.md

\section{Adaptation and merging}
\label{sec:adaptation-and-merging}
\label{sec:adaptation}

The previous sections (\S7--\S11) drove parameter updates with gradients
on data: SFT cross-entropy, RLHF/DPO preference signals, RLVR's
verifiable rewards. This section closes the pipeline with a different
class of operation altogether — \emph{weight-space} updates whose unit
of action is the parameter vector itself, not a per-token loss.
Parameter-efficient fine-tuning (PEFT) writes a low-rank patch onto a
frozen base; model merging averages, signs, and trims weight vectors
from fine-tunes that share an ancestor. Both rest on the same
empirical phenomenon — fine-tunes from a shared base land in a single
connected loss basin — and both are now production-load-bearing: every
open-weight LLM since LLaMA-2 ships with a LoRA adapter ecosystem, and
the strongest open multi-task models in 2024--2026 are merged from
specialists rather than jointly trained.

\subsection{LoRA: low-rank adaptation}
\label{sec:adapt-lora}

\citet{hu2021-lora} parameterize the update to a frozen pretrained
linear-layer weight $W_0 \in \R^{d \times k}$ as a low-rank
decomposition\footnote{KB excerpt:
\texttt{kb/excerpts/hu2021-lora.md\#sec-3}.}:
\begin{equation}
W = W_0 + \Delta W = W_0 + B A,
\quad
B \in \R^{d \times r},\ A \in \R^{r \times k},\ r \ll \min(d, k)
\label{eq:lora-decomp}
\end{equation}
\citep[\S 3]{hu2021-lora}. $W_0$ is frozen; only $A$ and $B$ are
trained. The forward pass at a hidden state $x$ is
\begin{equation}
h = W_0 x + \Delta W x = W_0 x + (\alpha / r)\, B A x,
\label{eq:lora-fwd}
\end{equation}
verbatim from \citet[\S 3]{hu2021-lora}, with scalar
$\alpha$ a constant in $r$ \citep[\S 3]{hu2021-lora}. The backward pass
flows gradients through $A$ and $B$ only:
\begin{equation}
\nabla_A \mathcal{L} = (\alpha/r)\, B^\top (\nabla_h \mathcal{L})\, x^\top,
\qquad
\nabla_B \mathcal{L} = (\alpha/r)\, (\nabla_h \mathcal{L})\, (A x)^\top,
\label{eq:lora-grad}
\end{equation}
so the optimizer state (Adam $m, v$) is allocated only over the rank-$r$
factors. The initialization is asymmetric: $A$ Gaussian, $B$ zero, so
that $\Delta W = 0$ at step 0 and the model behaves identically to the
base \citep[\S 3]{hu2021-lora}.

The trainable-parameter count drops from $\sum_\ell d_\ell k_\ell$ to
$\sum_\ell r(d_\ell + k_\ell)$. With LoRA applied to all attention
$Q, K, V, O$ projections and rank $r = 4$, GPT-3 175B's trainable count
falls roughly $10^{4}\times$ \citep[\S 3]{hu2021-lora}, with GPU memory
for fine-tuning down $\sim\!3\times$ \citep[abstract]{hu2021-lora}. The
empirical fact that \emph{rank $r = 4$ is sufficient for adapting the
weight matrices in the GPT-3 175B model} is reported verbatim in
\citet[\S 4]{hu2021-lora}, alongside the design choice that adapting
$Q$ and $V$ together gives the strongest task transfer.

The deployment-time mechanic is what separates LoRA from earlier
adapter families \citep{houlsby2019-adapters}\deepencite{houlsby2019-adapters
§3}. Adapter modules add bottleneck layers that incur inference
latency at every forward pass; LoRA computes
\begin{equation}
W_{\text{merged}} = W_0 + (\alpha / r)\, B A
\label{eq:lora-merge}
\end{equation}
\emph{once} at deployment \citep[\S 5]{hu2021-lora}, and the deployed
model has identical inference cost to the base. Different LoRAs can be
loaded and unloaded by adding or subtracting $(\alpha/r) B A$ — the
inference-time mechanic that S-LoRA and Punica multi-tenant systems
exploit (cross-ref Paper~1's serving discussion).

\subsection{QLoRA: quantize the base, train the adapter}
\label{sec:adapt-qlora}

\citet{dettmers2023-qlora} extend LoRA by quantizing the frozen base
$W_0$ to four bits while keeping the LoRA factors in BF16. The headline
\citep[abstract]{dettmers2023-qlora}\footnote{KB excerpt:
\texttt{kb/excerpts/dettmers2023-qlora.md\#abstract}.} is a 65B-parameter
fine-tune fitting on a single 48~GB GPU at full 16-bit task
performance. Three innovations carry the result:

\paragraph{NF4 — quantile-based 4-bit format.} The
\emph{NormalFloat-4} (NF4) format stores 16 bin centers chosen as the
quantiles of a standard normal distribution, then assigns each weight
to its nearest bin. Pretrained weights are empirically near-Gaussian
per block \citep[\S 3]{dettmers2023-qlora}\footnote{KB excerpt:
\texttt{kb/excerpts/dettmers2023-qlora.md\#sec-3-nf4}.}, so NF4's bin
spacing is information-theoretically near-optimal in this regime.
Concretely: bin centers are
\begin{equation}
q_i = \frac{1}{2}\!\left[Q_{\mathcal{N}(0,1)}\!\left(\tfrac{i}{16} + \tfrac{1}{2 \cdot 16}\right) + Q_{\mathcal{N}(0,1)}\!\left(\tfrac{i+1}{16} + \tfrac{1}{2 \cdot 16}\right)\right],
\quad i \in \{0,\ldots,15\},
\label{eq:nf4-bins}
\end{equation}
the midpoints of equal-mass quantiles of $\mathcal{N}(0,1)$
\citep[\S 3]{dettmers2023-qlora}, then re-scaled per block. Per-block
(block size 64) FP32 scales $Z$ are stored alongside, and the forward
GEMM dequantizes just-in-time:
\begin{equation}
y = \mathrm{dequant}(W_0^{\text{NF4}})\, x + (\alpha / r)\, B (A x),
\quad
\mathrm{dequant}(W_0^{\text{NF4}}) = Z \cdot \mathrm{NF4\_decode}(W_0^{\text{NF4}})
\label{eq:qlora-fwd}
\end{equation}
\citep[\S 3]{dettmers2023-qlora}, with $\mathrm{NF4\_decode}$ a 16-entry
lookup. The GEMM itself runs in BF16.

\paragraph{Double quantization.} The per-block FP32 scales are
themselves quantized to 8-bit with a second-level scale, dropping the
average overhead from $\sim\!0.5$~bits/parameter to $\sim\!0.37$~bits
\citep[\S 3]{dettmers2023-qlora}\footnote{KB excerpt:
\texttt{kb/excerpts/dettmers2023-qlora.md\#sec-3-doublequant}.}. At
65B parameters the saving is $\sim\!1$~GB.

\paragraph{Paged optimizers.} Optimizer state for the LoRA factors is
allocated in NVIDIA unified-memory pages that swap to CPU under GPU
pressure \citep[\S 3]{dettmers2023-qlora}\footnote{KB excerpt:
\texttt{kb/excerpts/dettmers2023-qlora.md\#sec-3-paged}.}. Long-sequence
batches that would OOM the 48~GB GPU instead spill to host RAM at a
small bandwidth cost.

QLoRA-tuned models recover within $\sim\!1\%$ of full BF16 fine-tunes
on 8 instruction datasets and multiple model scales
\citep[\S 4]{dettmers2023-qlora}, including 33B and 65B scales that
would otherwise be infeasible on commodity hardware.

\subsection{Task arithmetic: weight-space steering}
\label{sec:adapt-taskvec}

\citet{ilharco2022-task-vectors} reframe a fine-tune as a vector in
weight space. Given a pretrained $\theta_0 \in \R^P$ and a fine-tune
$\theta_T$ on task $T$, the \emph{task vector} is\footnote{KB excerpt:
\texttt{kb/excerpts/ilharco2022-task-vectors.md\#sec-2}.}
\begin{equation}
\tau_T = \theta_T - \theta_0 \in \R^P
\label{eq:task-vector}
\end{equation}
\citep[\S 2]{ilharco2022-task-vectors}. Three arithmetic operations on
task vectors are well-defined:
\begin{align}
\theta &= \theta_0 + \alpha\, \tau_T && \text{(addition; $\alpha = 1$ recovers $\theta_T$)}, \label{eq:tv-add}\\
\theta &= \theta_0 - \alpha\, \tau_T && \text{(negation; unlearns task $T$)}, \label{eq:tv-neg}\\
\theta &= \theta_0 + \sum_i \alpha_i\, \tau_{T_i} && \text{(combination; multi-task)} \label{eq:tv-combine}
\end{align}
\citep[\S 2]{ilharco2022-task-vectors}. The empirical claims
\citep[\S 3]{ilharco2022-task-vectors} are that negating a
toxicity-task vector reduces toxic generation while preserving general
benchmarks, that multi-task addition can outperform naive multi-task
fine-tuning, and that analogy compositions $\tau_A - \tau_B + \tau_C$
produce the expected behavioral shifts.

The mechanism, per \citet[\S 4]{ilharco2022-task-vectors}, is the
\emph{linearity of fine-tuning in weight space}: small movements in
the basin around $\theta_0$ produce small behavioral changes that sum
approximately linearly. We return to that claim — and its load-bearing
assumption of basin-sharing — in \cref{sec:adapt-lmc}.

\subsection{Model soups: averaging fine-tunes}
\label{sec:adapt-soups}

\citet{wortsman2022-model-soups}\footnote{KB excerpt:
\texttt{kb/excerpts/wortsman2022-model-soups.md\#sec-3}.} address the
hyperparameter-sweep waste problem. The conventional recipe trains $K$
fine-tunes with different hyperparameters, picks the best on
validation, and discards the rest \citep[\S 1]{wortsman2022-model-soups}.
The soup recipe averages them. The \emph{uniform soup} of $K$ fine-tunes
$\{\theta_k\}_{k=1}^K$ from the same base $\theta_0$ is
\begin{equation}
\theta_{\text{uniform}} = \frac{1}{K}\sum_{k=1}^K \theta_k
\label{eq:soup-uniform}
\end{equation}
\citep[\S 3]{wortsman2022-model-soups}. The \emph{greedy soup} sorts
checkpoints by validation accuracy and includes each $\theta_k$ in a
running average $\bar\theta$ only if doing so does not decrease
validation accuracy:
\begin{equation}
\bar\theta_k = \frac{|S_{k-1}|\, \bar\theta_{k-1} + \theta_k}{|S_{k-1}| + 1}
\quad
\text{if } \mathrm{ValAcc}(\bar\theta_k) \geq \mathrm{ValAcc}(\bar\theta_{k-1}),
\label{eq:soup-greedy}
\end{equation}
else skip \citep[\S 3]{wortsman2022-model-soups}. Both soups produce a
single inference-time model with no ensemble cost. The headline result
\citep[\S 1, \S 4]{wortsman2022-model-soups}: averaging fine-tunes from
a shared base \emph{often improves accuracy and robustness} over the
best single fine-tune, with no inference-time penalty, on CLIP, ALIGN,
and ViT-G fine-tuning sweeps. The condition that makes this work — the
fine-tunes lying in a single low-error basin — is the linear-mode-
connectivity assumption taken up in \cref{sec:adapt-lmc}.

\subsection{TIES-Merging: trim, elect sign, disjoint-merge}
\label{sec:adapt-ties}

Naive task-vector summation suffers \emph{interference}
\citep[abstract]{yadav2023-ties-merging}\footnote{KB excerpt:
\texttt{kb/excerpts/yadav2023-ties-merging.md\#sec-3}.}: when two
fine-tunes update the same coordinate in opposite directions, plain
addition cancels both. \citet{yadav2023-ties-merging} introduce a
three-step procedure that removes redundant low-magnitude updates and
resolves sign disagreement before averaging.

Given task vectors $\{\tau_i\}_{i=1}^K$ from a shared base $\theta_0$:

\paragraph{Step 1 --- Trim.} Keep the top-$k\%$ entries of each task
vector by magnitude, zero out the rest:
\begin{equation}
\tilde\tau_{i,j} =
\begin{cases}
\tau_{i,j} & \text{if } |\tau_{i,j}| \in \text{top-}k\% \text{ of } |\tau_i|,\\
0 & \text{otherwise}
\end{cases}
\label{eq:ties-trim}
\end{equation}
\citep[\S 3, step 1]{yadav2023-ties-merging}. Typical $k = 20\%$; the
trim removes \emph{interference due to redundant parameter values}
\citep[abstract]{yadav2023-ties-merging}.

\paragraph{Step 2 --- Elect sign.} For each coordinate $j$, take the
sign with the largest total magnitude across tasks:
\begin{equation}
s_j = \mathrm{sign}\!\left(\sum_{i=1}^K \tilde\tau_{i,j}\right)
\label{eq:ties-sign}
\end{equation}
\citep[\S 3, step 2]{yadav2023-ties-merging}. This is the
dominant-direction sign, weighted by magnitude — a magnitude-weighted
majority vote.

\paragraph{Step 3 --- Disjoint merge.} For each coordinate, average
only the task vectors whose sign matches the elected $s_j$:
\begin{equation}
\bar\tau_j = \frac{1}{|\mathcal{A}_j|}\sum_{i \in \mathcal{A}_j} \tilde\tau_{i,j},
\quad
\mathcal{A}_j = \{i : \mathrm{sign}(\tilde\tau_{i,j}) = s_j\}
\label{eq:ties-merge}
\end{equation}
\citep[\S 3, step 3]{yadav2023-ties-merging}. The merged model is
$\theta_0 + \lambda\, \bar\tau$ with scaling $\lambda$, typically $1$ or
tuned on a held-out set. Across vision and NLP benchmarks with up to
11 tasks merged, \citet[\S 4]{yadav2023-ties-merging} report the
ordering TIES $>$ Task Arithmetic $>$ Uniform Soup on multi-task
accuracy, with sign-conflict resolution contributing the largest gain
and trimming a smaller but consistent improvement.

\begin{intuition}
Task vectors are weight-space directions, and their arithmetic is
exactly that of vectors in $\R^P$. If $\tau_{\text{summarize-en}}$ is
the fine-tune-direction for English summarization and
$\tau_{\text{translate-fr}}$ is the fine-tune-direction for an
English-to-French translator, then
$\theta_0 + \tau_{\text{summarize-en}} + \tau_{\text{translate-fr}}$
empirically produces a French-summarization model
\citep[\S 3]{ilharco2022-task-vectors} — a behavioral analogy that
returns to canonical math via \cref{eq:tv-combine}: a positive linear
combination of task vectors lands at a weight that performs both
constituent tasks, modulo the interference TIES is designed to handle.
The mechanism TIES adds is that when the two task vectors disagree on
the sign of coordinate $j$ — push it $+$ for summarization and $-$ for
translation — \cref{eq:ties-merge} keeps only the dominant-magnitude
sign, so the merge does not silently zero out load-bearing
coordinates.
\end{intuition}

\subsection{Linear mode connectivity: why merging works}
\label{sec:adapt-lmc}

The shared substrate of \cref{sec:adapt-soups,sec:adapt-taskvec,sec:adapt-ties}
is the empirical claim that fine-tunes of a shared pretrained base land
in a \emph{connected} low-loss basin: the line segment
\begin{equation}
\theta(t) = \theta_0 + t\, (\theta_f - \theta_0),
\quad t \in [0, 1]
\label{eq:lmc-line}
\end{equation}
stays in low-loss territory. \citet[\S 3]{wortsman2022-model-soups}
provide the empirical evidence — loss is approximately monotone-low
along \cref{eq:lmc-line} for fine-tunes from a shared CLIP/ALIGN/ViT-G
base, and \emph{logit ensembling and weight averaging produce similar
accuracy when the loss landscape is locally flat}
\citep[\S 4]{wortsman2022-model-soups} — but they do not prove sufficient
conditions. The phenomenon is referred to in the broader literature as
\emph{linear mode connectivity} (LMC), with \citet{frankle2020-lmc}\deepencite{frankle2020-lmc
§3} the canonical reference\footnote{LMC was introduced for SGD-trained
networks at the same initialization; the fine-tuning case is the
operationally relevant specialization for adaptation.}; the
\citet[\S 4]{ilharco2022-task-vectors} \emph{linearity-of-fine-tuning}
attribution refers to the same phenomenon at the task-vector level.

The condition is empirically tight in a load-bearing way. LMC holds for
SFT-from-the-same-base; it holds partly for RLHF-from-the-same-base
(see \S8); it breaks for from-scratch models with different
initializations, where averaging produces a near-random model. This
explains the \emph{common-ancestor} requirement of every merge method:
LoRA's adapter is a delta on a fixed $W_0$, soup recipes assume the
$\theta_k$ share $\theta_0$, task arithmetic encodes the assumption in
\cref{eq:task-vector}'s subtraction, and TIES requires the ancestor
explicitly. Merging models with different pretrained bases is not
supported by any of these methods, because the assumed basin
isn't there.

\begin{contradiction}
The merge literature reports two patterns that do not yet reconcile.
On controlled multi-task benchmarks, merged models often exceed any
single ingredient \citep[\S 4]{yadav2023-ties-merging,
wortsman2022-model-soups}. On real-world deployments — the open-weights
LLM merging community at frontier scale — merged models routinely
under-perform the strongest specialist on its own task, with merge wins
concentrated in multi-task generalization rather than peak-task
performance. The per-paper benchmarks are largely single-source: TIES
\citep[\S 4]{yadav2023-ties-merging} and Soups
\citep[\S 4]{wortsman2022-model-soups} sweep their own task suites, with
no third-party head-to-head at frontier-LLM scale, and no canonical
theory predicts when merging produces interference vs.\ the
basin-centering effect predicted by LMC. The
\emph{merged-beats-best-ingredient} claim should be read as
\emph{task-set-conditional}: it holds on multi-task evaluations of
specialists, and is not yet established at the single-task ceiling.
\deepencite{yadav2023-ties-merging §6}
\end{contradiction}

\subsection{Forward link}

The pipeline closes here. The four sections following SFT (\S7) ---
RLHF (\S8), DPO (\S9), RLAIF/CAI (\S10), RLVR/GRPO (\S11) --- moved the
policy in data-driven directions; \cref{sec:adaptation-and-merging}
makes the dual move in weight space, with PEFT
(\cref{sec:adapt-lora,sec:adapt-qlora}) writing low-rank patches onto a
frozen base and merging
(\cref{sec:adapt-taskvec,sec:adapt-soups,sec:adapt-ties}) recombining
fine-tunes that share an ancestor. Both rest on the LMC substrate of
\cref{sec:adapt-lmc}, the speculative theoretical claim that holds the
empirical ceiling. The next section
(\cref{sec:frontier-snapshot})\deepencite{paper2-§13-cross-ref} fixes
the frontier-2026 picture: a single table over disclosed pre-training,
optimization, distributed, precision, SFT, preference-RL,
reasoning-RL, and merge choices for the open-weights frontier
(DeepSeek-V3/R1, Llama-3/4, Qwen-2.5/3, Phi-4, Tülu-3, OLMo-2, Gemma-3,
Mistral) and the partial disclosures from the closed labs. The
question \cref{sec:adaptation-and-merging} leaves open --- whether
merging will be a load-bearing primitive of the frontier pipeline or
a community add-on layered on top of single-base specialists --- is one
of the open boundaries \cref{sec:contradictions}\deepencite{paper2-§14-cross-ref}
takes up.
